% ------------------------------------------------------------
% Section 1 -- Introduction
% ------------------------------------------------------------
\section{Introduction}
\label{sec:intro}

Low-dose CT reduces radiation exposure at the cost of elevated image noise, and deep learning has become the dominant approach to restoring diagnostic quality. Since RED-CNN~\cite{chen2017}, a long line of architectures has been evaluated almost exclusively with pixel-level metrics --- PSNR, SSIM, VIF --- computed against normal-dose references. Yet a recent standardized benchmark found that on precisely these metrics, ``most methods perform statistically similar and improvements over the past six years have been marginal at best''~\cite{eulig2024}. We take this finding seriously, not as a reason to chase further decimal points, but as evidence that the classical evaluation axis is saturated --- and that what these metrics measure is not what remains broken.

What remains broken is \emph{physics} --- and the root cause is the training objective, not any particular architecture. Appearance-oriented losses (pixel MSE and its variants, but equally SSIM-, perceptual-, and adversarially-driven objectives) reward outputs that \emph{look} close to the reference while remaining agnostic to the statistics of the noise being removed and to the absolute calibration of the values being reconstructed. Any denoiser trained this way --- CNN, GAN, transformer, or diffusion-based --- is therefore free, and often incentivized, to suppress noise unevenly across spatial frequencies and to shift tissue attenuation wherever doing so lowers the average pixel error. We audit this failure mode on a benchmark-faithful, MSE-trained RED-CNN, chosen as a representative and standardized testbed rather than as a uniquely flawed model: despite being the strongest configuration by SSIM in our study, it under-removes low-frequency noise by nearly an order of magnitude in the abdomen, leaves 26\% of abdominal noise magnitude in the image, reproduces the true noise texture with a spectral-shape correlation of only 0.51, and introduces systematic, tissue-dependent HU calibration errors --- up to $-53$\hu{} in chest bone --- that no pixel-average metric registers.

Why should these two physical axes matter clinically? CT is a quantitative modality: Hounsfield units are calibrated physical measurements, and clinical decisions read them directly --- adrenal-adenoma and fatty-liver characterization rely on attenuation thresholds of a few tens of HU, coronary-calcium scoring counts voxels above fixed HU cutoffs, radiotherapy planning converts HU to electron density, and radiomics and longitudinal follow-up assume HU stability across scans. A denoiser that silently shifts bone by $-53$\hu{}, or fat by tens of HU, can flip such decisions while improving PSNR. The noise power spectrum is equally consequential: radiologists are trained on the characteristic texture of CT noise, and a spectrally distorted residual produces the over-smoothed, ``plastic'' appearance that can erode diagnostic confidence; model-observer studies show that noise texture --- not merely noise magnitude --- governs low-contrast lesion detectability~\cite{solomon2012,solomon2020}, and downstream CAD and radiomics pipelines implicitly assume realistic noise statistics. An LDCT denoiser can therefore be simultaneously excellent by PSNR/SSIM and physically unfaithful in ways that may carry direct clinical consequences.

A natural response is loss-level: penalize NPS mismatch during training, as explored in recent work~\cite{li2025,zhang2025}. We confirm that a radial log-NPS loss paired with a per-tissue HU-bin bias loss recovers much of the fidelity at near-zero classical cost --- but we also identify a \textbf{ceiling}: doubling the NPS weight without any architectural change makes every metric worse, physics included. The gradient signal exists, but the tested trunk, under the matched training setup, cannot exploit the stronger spectral incentive without degradation --- it lacks a dedicated low-dimensional mechanism to express targeted spectral corrections.

We therefore propose an architectural mechanism: a \textbf{DC-preserving, anatomy-adaptive radial spectral head} (Fig.~\ref{fig:architecture}). The head leaves the trunk untouched and operates on the removed-noise residual in the Fourier domain, multiplying it by a learnable radial gain curve with structural guarantees that hold per image by construction --- exact DC preservation of the residual correction, frozen near-DC gains that prevent the head from altering large-area regional HU means, an isotropic 32-knot radial parameterization that introduces no new directional frequency content, and adaptive per-image offsets bounded by $\tanh$. A ${\approx}5$k-parameter conditioning encoder makes the curve anatomy-adaptive, and a batch-centering penalty --- motivated by a failure mode we measure and diagnose --- prevents the adaptive path from degenerating into a trunk-absorbable static shift. At strictly matched budgets and loss weights, the head wins every physics metric at equal classical quality, and its learned gain curves separate chest from abdomen by ${\approx}4\times$ the within-anatomy spread while every guarantee holds.

\paragraph{Contributions.}
\begin{enumerate}
  \item \textbf{A physics audit of a classically strong baseline}, showing that appearance-driven training produces spectrally and calibrationally unfaithful denoising invisible to PSNR/SSIM --- a failure mode generic to image-appearance objectives, quantified here on RED-CNN as a representative testbed --- and an evaluation protocol (matched budgets, per-anatomy reporting, physics metrics) that exposes it.
  \item \textbf{A DC-preserving, anatomy-adaptive spectral head with a trunk-agnostic interface} (${\approx}5$k parameters, one FFT pair at inference) whose per-image guarantees hold on any trunk by construction --- while its quantitative benefit is measured per trunk (Sec.~\ref{sec:crosstrunk}) --- and a centering constraint that we show is necessary to keep adaptive capacity anatomy-discriminative rather than statically degenerate.
  \item \textbf{An attribution study} at matched budgets and matched loss weights demonstrating that the head --- not the losses --- accounts for the best physical fidelity in the study (NPS error $-12\%$ overall and $-19\%$ in the chest, chest biases substantially reduced, at equal image quality), including a loss-ceiling experiment showing losses alone cannot exploit stronger physics incentives.
  \item \textbf{Quantified adaptivity}: to our knowledge the first LDCT module whose anatomy-adaptive behavior is itself measured (separation vs.\ within-anatomy spread) rather than asserted, with directly interpretable gain-curve visualizations.
\end{enumerate}

The remainder of the paper is organized as follows: Sec.~\ref{sec:related} reviews related work; Sec.~\ref{sec:method} describes the method; Sec.~\ref{sec:setup} details data, protocol, and metrics; Sec.~\ref{sec:results} presents results and ablations; Sec.~\ref{sec:discussion} discusses findings and limitations.
